\begin{claim}{D.2.23}{Proved}{The actual canonical Browder square vanishes}
For the class $\chi\in\pi_4F$ of \textup{D.2.21},
\[
 \lambda_1(\chi,\chi)=0\in\pi_9F,
 \qquad [\widehat\chi,\widehat\chi]=0\in\pi_{11}Y.
 \tag{D.2.23a}
\]
Consequently the fixed canonical summand of \textup{D.2.21}
extends through relative dimension ten. This assertion does not
extend the separate correction or the combined map $G_9$.
\end{claim}
\begin{proof}
The integral Hurewicz calculation and complete kernel computation
of \textup{D.2.22} put the actual square in
$\{0,\rho\nu_6\}$, with $\rho\nu_6\ne0$.
We distinguish these two actual maps by a cohomology relation
and their ordinary mapping cones.

Use the homology bases of \textup{D.2.22}, and write
$e_6^*,J_1^*$ for their dual cohomology vectors. The entire
eighteen-row table \textup{(D.2.22a6)} gives
\[
 \Psi=\operatorname{Sq}^4e_6^*
              +\operatorname{Sq}^2J_1^*=0
                 \quad\text{in }H^{10}(F;\FF_2).
 \tag{D.2.23b}
\]
Indeed each summand evaluates to one precisely on $B_4,B_5$
and to zero on the other sixteen basis vectors. In particular
the fourth squares of $D_2,Q_{10},P_p$ are respectively
$d_6,f_6,f_6$, with zero $e_6$ coefficient.

For any $f:S^9\to F$ with zero mod-two Hurewicz image,
the two lower cohomology classes extend uniquely to
$C_f=F\cup_f e^{10}$. The expression in
\textup{(D.2.23b)} is zero on the old $H_{10}F$ and hence
has a well-defined value on a new top homology class.
This value is independent of adding an old homology class.
For $f=\rho\nu_6$ the ordinary map of cones from
$S^6\cup_{\nu_6}e^{10}$ gives
\[
 \operatorname{Sq}_*^4(v)=h(\rho)=e_6+f_6,
 \qquad \operatorname{Sq}_*^2(v)=0.
 \tag{D.2.23c}
\]
The first is the Hopf-invariant-one detector of $\nu$,
and the second follows from the zero degree-eight homology
of this two-cell source. Thus the value of $\Psi$ is one
on this cone. Its value on the null map is zero.

Now form the cone of the actual Browder square. Since $F$ is
an actual double-loop space, $\chi$ has a free $E_2$ extension
$C_2S^4\to F$. In $C_2S^4\to C_3S^4$ the only relative
cell through dimension ten is the binary ten-cell. Its
attaching map is the Browder square, with the normalization
proved in \textup{D.2.21}. Push out this relative ten-stage
along the free extension. The result is the ordinary
homotopy cofiber
\[
 C_{\lambda_1(\chi,\chi)}=F\cup_{\lambda_1(\chi,\chi)}e^{10}.
 \tag{D.2.23d}
\]
No $E_3$ structure on $F$ is assumed in this pushout.

Let $x\in H_4S^4$ be the bottom class. In the source the
absolute class $v=Q^6x$ has coefficient one on the relative
ten-cell. This follows either from its binary configuration
representative or the free three-fold loop homology basis
\cite[Chapter III, Theorems 3.1--3.2]{I_CohenLadaMay1976}.
There is no degree-ten class in $H_*C_2S^4$, and the mod-two
Browder attaching boundary is zero, so this top coefficient
is unambiguous. The source comparison $C_3S^4\to QS^4$ is
ten-connected. It is therefore injective on $H_8$, and the
ordinary stable Nishida formula gives, already in $C_3S^4$,
\[
 \operatorname{Sq}_*^2(Q^6x)
 =\binom42 Q^4x+Q^5\operatorname{Sq}_*^1x=0,
 \qquad \operatorname{Sq}_*^4(Q^6x)=0.
 \tag{D.2.23e}
\]
Here $\binom42=0$ modulo two and every positive square of
$x$ is zero. The second equality also follows from $H_6C_3S^4=0$.
These operations are unchanged in the relative ten-stage,
whose comparison with the full source is an isomorphism
in degrees six and eight. The Nishida input is
\cite[Chapter III, Theorem 1.3]{I_CohenLadaMay1976}; the
comparison with the infinite-loop source avoids any
unstated top-operation correction.

The pushout map takes $v$ to a class with coefficient one
on the new cone cell and with both displayed squares zero,
by ordinary Steenrod naturality. Therefore the value of
$\Psi$ on this cone is zero. Since the only two possibilities
were $0$ and $\rho\nu_6$, and the latter has value one by
\textup{(D.2.23c)}, the actual square is zero. The extension
conclusion follows from \textup{D.2.21}.
\end{proof}
